%%%%%%%%%%%%%%%%%%%%%%%%%%%%%%%%%
%
%       EXERCÍCIO
%
%%%%%%%%%%%%%%%%%%%%%%%%%%%%%%%%%

\definevalorquestao

\begin{exercicioBanco}[\valorquestao]
O acréscimo de tecnologias no sistema produtivo industrial tem por objetivo reduzir custos e aumentar a produtividade. No primeiro ano de funcionamento, uma indústria fabricou \(8.000\) unidades de um determinado produto. No ano seguinte, investiu em tecnologia adquirindo novas máquinas e aumentou a produção em \(50\%\). Estima-se que esse aumento percentual se repita nos próximos anos, garantindo um crescimento anual de \(50\%\). Considere \(P\) a quantidade anual de produtos fabricados no ano \(t\) de funcionamento da indústria. Se a estimativa for alcançada, qual é a expressão que determina o número de unidades produzidas \(P\) em função de \(t\), para \(t \ge 1\)?

% Define as alternativas
\newcommand{\alternativas}{%
    \begin{center}
        \begin{tabularx}{\textwidth}{XXX}
            (a) \(P(t) = 0{,}5^{t-1} + 8000\). &
            (b) \(P(t) = 50 \cdot (0{,}5)^{t-1} + 8000\). &
            (c) \(P(t) = 4000 \cdot 2^{t-1} + 8000\). \\[5pt]
            (d) \(P(t) = 8000 \cdot (0{,}5)^{t-1}\). &
            (e) \(P(t) = 8000 \cdot (1{,}5)^{t-1}\).
        \end{tabularx}
    \end{center}
}

% Define a resposta correta
\newcommand{\resposta}{E}

% Lógica condicional para exibição
\ifdefstring{\atividade}{lista}{%
    \alternativas
    \vspace{0.5em}
    
    \noindent\textbf{Resposta:} letra \textbf{\resposta}.
}{%
    \ifdefstring{\modo}{objetivo}{%
        \alternativas
    }{%
        % modo = discursiva → não mostra alternativas
    }
}
\end{exercicioBanco}

